% Ad Familiares 12.25 --- headnote
% ad-familiares-12-25, 19 March 43 BC, Rome

\textit{Cicero to Q. Cornificius, governor of Old Africa, from Rome
on 19 March 43 BC --- Perseus dateline \textit{Scr. Romae xiv K. Apr.
a. 711 (43)}. The letter is one of the more important in the
Cornificius file: a retrospective of Cicero's whole campaign against
Antonius, written in the mood after Hirtius has marched out for Mutina
and the political wind in Rome has turned squarely against the
Antonian camp. Cicero opens with a piece of Senate news --- on the
Quinquatrus, the feast of Minerva, a complimentary decree was passed
in Cornificius's honour, to the discomfiture of two of his enemies,
the ``Minotaur'' Calvisius and Taurus --- and ties it elegantly to
the day's omen: Minerva's statue, dislodged by a whirlwind, was on
the same day voted restored.}

\textit{The body of the letter is the autobiography of the Philippics
in miniature. Cicero rehearses his 20 December 44 BC motion in the
Senate (``the thirteenth day before the Kalends of January'') as the
foundation laying of the new state; his goading of Antonius into the
fury that drove him out of Rome; the storm that turned his own ship
back from his Greek flight (\emph{the Etesian winds, like good
citizens, refused to escort me as I deserted the state}); and the
celebrated and unflattering claim that the young Caesar Octavianus
caught Antonius ``belching and retching'' in his ``nets.'' He closes
with a Terence quotation (\emph{Andria} 189) and a sustained
ship-of-state metaphor: there is now one ship for all good men, and
Cornificius is invited to come aboard --- and indeed to the stern.}
